\documentclass{situationist-preprint}

\title{Time Without Change}
\msid{TSM-17-a}
\msfamily{TSM-17}
\msdiscipline{Time, Space and Modality}
\mstype{Principal manuscript}
\mskeywords{time without change; temporal vacuum; freeze worlds; Shoemaker; relationism; simplicity}

\begin{document}
\maketitle

\begin{abstract}
Aristotle held that there is no time without change, and the claim has been repeated by relationists, verificationists and McTaggart. We argue that time without change is possible and could be known to occur. The argument is Shoemaker's. In a world divided into three regions, each region freezes at regular intervals, every third, fourth and fifth year respectively, for a year at a time; the inhabitants of the other regions observe each local freeze, and its inhabitants learn of it on their return. The simplest law that fits their observations implies that every sixtieth year all three regions freeze together, so that for a year nothing anywhere changes. The inhabitants would be justified in believing this, and if the law is true they would know it. We state the argument with explicit premises, show that the rival hypothesis on which joint freezes never occur is an unmotivated equivalent of the kind that does not threaten justification, and answer the objection that nothing could end a total freeze. The result fits a tenseless view of time: the B-series is an order of times, and it needs no change to be populated.
\end{abstract}

\section{The question}

\citet[IV.11, 218b21--219a10]{AristotlePhys} observed that when we notice no change in our minds, we do not notice that time has passed, and concluded that time does not exist without change. \citet[1.2.3]{HumeTreatise} held that the idea of time is derived from the succession of our perceptions, so that the idea of a time in which nothing changes is a fiction. Relationists, following Leibniz \citep{Alexander1956}, take time to be nothing over and above the order of successive events, so that a stretch of time in which nothing happens would be a stretch of nothing. \citet{McTaggart1908} made the dependence of time on change the first premise of his argument for the unreality of time.

\citet{Shoemaker1969} answered with a thought experiment. Its target is the claim that nobody could have reason to believe that a period without change had occurred; it describes a world whose inhabitants would have such reason. \citet{NewtonSmith1980} accepted that the description is coherent but argued that the evidence underdetermines the hypothesis of total freezes and their duration. We defend Shoemaker's conclusion in full: time without change is possible and could be known to occur.

\section{Framework}

\begin{definition}[Change]\label{d:change}
A thing \emph{changes} if it has a qualitative property at one time that it lacks at another. Mere variation in temporal relations, such as an object's coming to be one year older, is not change in this sense.
\end{definition}

\begin{definition}[Freeze]\label{d:freeze}
A \emph{local freeze} is an interval during which nothing in a region changes. A \emph{total freeze} is an interval during which nothing in the universe changes.
\end{definition}

\begin{definition}[Freeze world]\label{d:world}
A \emph{freeze world} is divided into regions $A$, $B$ and $C$, whose inhabitants can observe one another. Region $A$ freezes for one year in every third year, $B$ in every fourth and $C$ in every fifth. A frozen region looks to outside observers as if time had stopped within it; when the freeze ends, things in it resume exactly as they would have had the freeze lasted an instant, and its inhabitants learn from others that a year has elapsed.
\end{definition}

\begin{assumption}\label{a:observable}
Local freezes are observable from outside the frozen region, and their periods can be measured by clocks that continue to run elsewhere.
\end{assumption}

\begin{principle}[Simplicity]\label{p:simple}
If two hypotheses fit the evidence equally well, and one states the regularities found in the evidence without exception while the other adds an exception the evidence does not motivate, belief in the first is justified over belief in the second.
\end{principle}

Principle~\ref{p:simple} is an instance of ordinary inductive practice, which EPI-03 argues is justified. EPI-04 argues that projectibility tracks naturalness; a regularity with an exception tailored to unobservable cases is an unnatural, grue-like predicate.

\section{The argument}

\begin{proposition}\label{p:possible}
A freeze world is possible, and in it a total freeze occurs every sixtieth year.
\end{proposition}

\begin{proof}
The description contains no contradiction: each local freeze is an interval during which the things in one region stand in the same qualitative state, and nothing in the concept of a region or of a time requires that some property of the region vary over every interval. Since $3$, $4$ and $5$ are pairwise coprime, the years in which all three regions freeze recur with period $60$, the least common multiple. In those years nothing in the universe changes.
\end{proof}

\begin{proposition}\label{p:local}
The inhabitants of a freeze world are justified in believing that each region freezes with its observed period.
\end{proposition}

\begin{proof}
By Assumption~\ref{a:observable}, each local freeze is observed by the inhabitants of the other two regions, who record its onset, end and recurrence. Over a long enough record each period is confirmed by many instances, which is ordinary inductive evidence for a regularity.
\end{proof}

\begin{theorem}\label{t:main}
Time without change is possible, and it could be known to occur.
\end{theorem}

\begin{proof}
Possibility is Proposition~\ref{p:possible}. For knowability: by Proposition~\ref{p:local}, the inhabitants are justified in believing the three periodic laws. Two hypotheses fit their observations. On the first, the three laws hold without exception, and so, by the arithmetic of Proposition~\ref{p:possible}, a total freeze occurs every sixtieth year. On the second, the laws hold except in the years in which they would coincide, in which, without any change in what is observed, no region freezes. No possible observation distinguishes them, since in a total freeze nobody observes anything; but the second adds an exception that nothing in the evidence motivates. By Principle~\ref{p:simple}, the inhabitants are justified in believing the first, and so in believing that total freezes occur. The inhabitants never perceive a total freeze, which is all that Aristotle's and Hume's observation concerns; they infer it from regularities they do observe. In a freeze world the belief is true, and it is held on the basis of the regularities that make it true, which is the kind of basis on which inductive knowledge of laws is had. So the inhabitants could know that time without change occurs.
\end{proof}

\section{Objections and replies}

\begin{objection}[Underdetermination]
The no-joint-freeze hypothesis is empirically equivalent to Shoemaker's, and so are hypotheses on which the joint freeze lasts a day, or a century. The evidence leaves the occurrence and the length of total freezes undetermined \citep{NewtonSmith1980}.
\end{objection}

\begin{reply}
Each rival is obtained by adding to the uniform laws a clause about precisely the cases that cannot be observed. SCI-04 argues that theories are underdetermined only by empirically equivalent rivals that lack independent motivation, and these rivals are of that kind. The length of a joint freeze is fixed by the same laws: each local freeze lasts a year, so a year in which all three regions freeze is a year without change. A hypothesis that assigns the joint freeze another length must assign the local freezes a different behavior in exactly the unobserved case, and is unmotivated in the same way.
\end{reply}

\begin{objection}[Nothing could end it]
If nothing changes during a total freeze, then the state of the world at its end is the same as at its beginning. Whatever caused the world to resume would have caused it to resume at once; so nothing could end a total freeze, and the world would stay frozen forever.
\end{objection}

\begin{reply}
The objection assumes that a cause must be a state that immediately precedes its effect. The laws of a freeze world relate a region's state to its state a year later: a region that enters a freeze in a certain state resumes after the elapse of one year. Such a law makes the elapsed interval a causal factor, and this is what \citet{Shoemaker1969} called action at a temporal distance. It is unusual, but it is not incoherent, and the inhabitants have evidence for it in every local freeze they observe ending. A world with such laws is the world described in Definition~\ref{d:world}, and nothing in it requires change during the freeze.
\end{reply}

\begin{objection}[Time is change]
Time just is the succession of changes. Where there is no change there is no succession, and so no time; what Shoemaker describes is at best a world whose clocks skip a year.
\end{objection}

\begin{reply}
This is the relationist premise that the argument puts in question, and it cannot be assumed against it. On the tenseless view defended in TSM-01, time is constituted by the B-series, an order of times by the relations of earlier and later. A total freeze is a stretch of that order at whose times the world's states are qualitatively alike. If things persist by having temporal parts at successive times, as PID-11 holds of persons, then a frozen thing has many temporal parts during a freeze, each qualitatively like the last; they are many and ordered, though not different in character. McTaggart's premise that time requires change was a premise about the change of A-properties, which TSM-01 rejects; the B-series needs no change to be populated.
\end{reply}

\section{What the argument does not establish}

The argument does not show that our world contains a total freeze, nor that it does not. Its conclusion is modal and epistemic: such a period is possible and could be known. It establishes a constraint on theories of time, namely that no theory on which time is identical with change is necessary, and so it tells against reductive relationism about time. It does not settle substantivalism about time, which would need the further claim that times are entities; TSM-06 defends the corresponding claim for space on separate grounds.

\section{Conclusion}

The claim that there is no time without change rests on the observation that we notice no time when we notice no change. Shoemaker's freeze world shows that what cannot be perceived can be inferred. Its inhabitants observe local freezes, infer the laws that govern them, and are justified by ordinary inductive standards in believing that every sixtieth year the whole world stands still. The description is coherent, the inference is sound, and if the laws hold the belief is knowledge. Time without change is possible and could be known to occur.

\begin{thebibliography}{99}\small
\bibitem[Alexander(1956)]{Alexander1956} Alexander, H.~G. (ed.) (1956). \emph{The Leibniz--Clarke Correspondence}. Manchester University Press.
\bibitem[Aristotle(Phys.)]{AristotlePhys} Aristotle. \emph{Physics}. Cited by book, chapter and Bekker page.
\bibitem[Hume(1739)]{HumeTreatise} Hume, D. (1739). \emph{A Treatise of Human Nature}. Cited by book, part and section.
\bibitem[McTaggart(1908)]{McTaggart1908} McTaggart, J.~M.~E. (1908). The unreality of time. \emph{Mind}, 17(68), 457--474.
\bibitem[Newton-Smith(1980)]{NewtonSmith1980} Newton-Smith, W.~H. (1980). \emph{The Structure of Time}. Routledge and Kegan Paul.
\bibitem[Shoemaker(1969)]{Shoemaker1969} Shoemaker, S. (1969). Time without change. \emph{The Journal of Philosophy}, 66(12), 363--381.
\end{thebibliography}

\end{document}
